\chapter{Diseño metodológico: \emph{point-in-time} y ausencia de \emph{lookahead}}
\label{cap:metodologia}

Si el capítulo anterior argumentaba que el rigor empieza en la base de datos, este desarrolla el
mecanismo concreto que lo garantiza: la construcción \emph{point-in-time}. La regla que gobierna
todo el sistema es sencilla de enunciar y difícil de respetar sin descuido: \textbf{una señal
fechada un día solo puede usar información observable ese día}. El capítulo formaliza la regla,
muestra por qué la aproximación ingenua ---aplicar un retardo fijo--- la incumple, describe el
algoritmo de observabilidad y explica por qué los contrastes de fuga temporal se escriben antes que
el código, no después.

\section{Notación}
\label{sec:notacion}

Se fija aquí la notación que el resto de la memoria reutiliza sin volver a definir:

\begin{itemize}
  \item \(\tsnap\): la \textbf{fecha de \emph{snapshot}}, el día en que se genera una señal y se
  reconstruye lo que era observable hasta ese instante.
  \item \(\tfiled\): la \textbf{fecha real de publicación} de un dato fundamental, tal como la
  registra SEC EDGAR (\texttt{filingDate}).
  \item \(\hlabel\): el \textbf{horizonte de la etiqueta}, el número de meses hacia el futuro sobre
  el que se mide el retorno que el modelo intenta ordenar.
\end{itemize}

La regla \emph{point-in-time} se escribe entonces como una condición de observabilidad: un
fundamental con fecha de publicación \(\tfiled\) solo puede intervenir en una señal de fecha
\(\tsnap\) si \(\tfiled \le \tsnap\). La etiqueta asociada a esa señal se realiza en
\(\tsnap + \hlabel\) y, por construcción, es información del futuro que nunca entra en el
entrenamiento hasta que se ha realizado (capítulo~\ref{cap:agentes}).

\section{Por qué el retardo fijo es incorrecto}

La tentación natural al no disponer de la fecha de publicación es aproximarla: suponer que todo
informe es público, pongamos, 45 días después de su cierre fiscal, y aplicar ese retardo constante
a todas las empresas. El enfoque es cómodo pero \textbf{metodológicamente incorrecto}, y el proyecto
lo descartó tras contrastarlo con las fechas reales de EDGAR.

El motivo es que el retardo real de publicación es muy variable entre empresas y entre informes. En
los datos del proyecto se documentaron casos en los que una compañía tardó \textbf{133 días} en
publicar (AT\&T) y un informe anual (10-K) de Apple tardó \textbf{88 días}. Un retardo fijo de 45
días trataría esos datos como públicos semanas antes de que lo fueran realmente, inyectando
\emph{lookahead bias} justo en las empresas que más tardan ---precisamente aquellas cuya demora
puede correlacionar con dificultades que el modelo no debería anticipar---. La conclusión fue
inequívoca: cualquier retardo fijo introduce fuga, y la única forma correcta de saber cuándo un dato
era observable es usar la fecha en que de verdad se publicó.

\section{El papel de \texttt{lag\_days}: margen de ejecución, no observabilidad}

Una consecuencia importante de lo anterior es una \textbf{redefinición} de un parámetro del sistema.
El parámetro \texttt{lag\_days} deja de significar ``retardo con el que un dato se vuelve
observable'' ---esa función la cumple ahora \(\tfiled\)--- y pasa a significar únicamente el
\textbf{margen de ejecución}: el número de días que simula el desfase entre que se dispondría de la
información y el día en que realmente se lanzaría el \emph{pipeline} para operar. Es una holgura
operativa, conservadora, que se suma a la observabilidad real; no la sustituye ni la aproxima.

\section{Algoritmo de observabilidad}

La construcción del panel \emph{point-in-time} (\texttt{module/data/dataset.py}) recorre las fechas de
\emph{snapshot} y, en cada una, reconstruye lo observable: para cada empresa perteneciente al índice
en esa fecha (capítulo~\ref{cap:datos}), toma su \textbf{último informe realmente publicado},
esto es, el informe con mayor \(\tfiled\) que cumpla \(\tfiled \le \tsnap\). El resultado es un panel
indexado por \((\text{ticker}, \tsnap)\) en el que cada fila contiene solo información que existía en
ese instante. No hay relleno hacia atrás de datos que aún no se habían publicado, ni se mezclan
magnitudes de periodicidades distintas en un mismo corte.

\section{Tres fugas silenciosas y por qué los tests van primero}

Durante la auditoría del \emph{pipeline} aparecieron tres defectos que no rompían ninguna prueba
existente pero producían números sutilmente falsos. Se documentan porque ilustran que, en un
proyecto cuyo eje es la ausencia de \emph{lookahead}, el peligro no son los errores ruidosos sino
los silenciosos:

\begin{itemize}
  \item \textbf{Crecimiento interanual por posición.} El cálculo del crecimiento interanual
  comparaba ``cuatro trimestres atrás'' contando por posición en la serie; cuando faltaba un
  trimestre ---algo frecuente--- comparaba contra una fecha arbitraria. Un caso real producía un
  crecimiento falso del orden del \(+71\,\%\) al comparar contra un dato de quince meses antes.

  \item \textbf{Etiquetas perdidas con el día de \emph{snapshot} desplazado.} La fecha de la
  etiqueta se derivaba por aritmética de calendario, que trata los fines de mes de forma distinta a
  la rejilla de fechas; con el día 31 se perdía cerca del 40\,\% de las etiquetas.

  \item \textbf{Mezcla de magnitudes anuales y trimestrales.} Un margen anual y el de su cuarto
  trimestre comparten fecha de cierre; un \emph{fallback} mal planteado confundía magnitudes de 12 y
  de 3 meses en el mismo corte.
\end{itemize}

Los tres se corrigieron acompañados de \textbf{tests de regresión} que fallan contra el código
anterior. La lección metodológica ---que este proyecto adopta como norma--- es que los contrastes
de fuga temporal y de separación entre entrenamiento y evaluación se escriben \textbf{antes} que la
funcionalidad que verifican: son la primera línea de defensa del resultado, no un añadido posterior.
